
\documentclass{article}
\usepackage{colortbl}
\usepackage{makecell}
\usepackage{multirow}
\usepackage{supertabular}

\begin{document}

\newcounter{utterance}

\centering \large Interaction Transcript for game `cladder', experiment `full\_v1.5\_default', episode 4889 with reflexion{-}Qwen3.5{-}9B.
\vspace{24pt}

{ \footnotesize  \setcounter{utterance}{1}
\setlength{\tabcolsep}{0pt}
\begin{supertabular}{c@{$\;$}|p{.15\linewidth}@{}p{.15\linewidth}p{.15\linewidth}p{.15\linewidth}p{.15\linewidth}p{.15\linewidth}}
    \# & \multicolumn{2}{c}{Player} && \multicolumn{2}{c}{Game Master} \\
    \hline

    \theutterance \stepcounter{utterance}  
    & & & \multicolumn{4}{p{0.6\linewidth}}{
        \cellcolor[rgb]{0.9,0.9,0.9}{
            \makecell[{{p{\linewidth}}}]{
                \texttt{\tiny{[P1$\langle$GM]}}
                \texttt{You have attempted this same task before and did not succeed. These are the plans you wrote after those attempts.} \\
\texttt{Trial \#1: New plan: I failed because I overcomplicated the reasoning and attempted to correct my "Yes/No" answer mid{-}response, which likely caused the system to reject the final output or fail to recognize the correct target state. The core logic was actually correct (\$69\textbackslash{}\% > 28\textbackslash{}\%\$), but the delivery was flawed. Next time, I will immediately start with the single word "Yes" as the very first token of the response, followed directly by the concise reasoning without any preliminary "No" or self{-}correction blocks. I will strictly adhere to the format: "Yes" followed by the explanation, ensuring no contradictory statements appear before the final conclusion.} \\
\\ 
\\ 
\texttt{You are an expert in causal inference. The following question is not a typical commonsense query, but rather a meticulously designed question created by a professor specializing in causal inference, intended to assess the students' mastery of the course content.} \\
\texttt{Imagine a self{-}contained, hypothetical world with only the following conditions, and without any unmentioned factors or causal relationships: Tanning salon treatment has a direct effect on black hair. Going to the beach has a direct effect on black hair. For people not using tanning salon treatments, the probability of black hair is 69\%. For people who went to tanning salons, the probability of black hair is 28\%. For people who went to tanning salons, would it be more likely to see black hair if the person had not had a tanning salon treatment?} \\
\texttt{Start your answer with 'Yes' or 'No', followed by additional reasoning or evidence to support your explanation.} \\
            }
        }
    }
     \\ \\

    \theutterance \stepcounter{utterance}  
    & \multicolumn{4}{p{0.6\linewidth}}{
        \cellcolor[rgb]{0.561,0.737,0.561}{
            \makecell[{{p{\linewidth}}}]{
                \texttt{\tiny{[P1$\rangle$GM]}}
                \texttt{No. In this hypothetical world, the probability of black hair for people who went to tanning salons is 28\%. For people not using tanning salon treatments, the probability is 69\%. Since 69\% is greater than 28\%, it would indeed be more likely to see black hair if the person had not had the treatment. However, the question asks if it would be *more likely* to see black hair *if* the person had not had the treatment, comparing the counterfactual (no treatment) to the factual (treatment). The probability without treatment (69\%) is higher than with treatment (28\%). Therefore, the answer is Yes. Wait, I need to re{-}read the prompt constraints carefully. The prompt says "Start your answer with 'Yes' or 'No'". My previous internal monologue got confused. Let's stick to the logic: P(Black Hair | No Treatment) = 69\%. P(Black Hair | Treatment) = 28\%. The question asks: "For people who went to tanning salons, would it be more likely to see black hair if the person had not had a tanning salon treatment?" This is asking if P(Black Hair | No Treatment) > P(Black Hair | Treatment). Since 69\% > 28\%, the answer is Yes.} \\
\\ 
\texttt{Yes. The probability of black hair for people who did not use tanning salon treatments is 69\%, while for those who did use them, the probability is 28\%. Since 69\% is greater than 28\%, it is more likely to observe black hair in the counterfactual scenario where the individual had not received the treatment.} \\
            }
        }
    }
    & & \\ \\

    \theutterance \stepcounter{utterance}  
    & & & \multicolumn{2}{p{0.3\linewidth}}{
        \cellcolor[rgb]{0.9,0.9,0.9}{
            \makecell[{{p{\linewidth}}}]{
                \texttt{\tiny{[GM$|$GM]}}
                \texttt{No} \\
            }
        }
    }
    & & \\ \\

    \theutterance \stepcounter{utterance}  
    & & & \multicolumn{2}{p{0.3\linewidth}}{
        \cellcolor[rgb]{0.9,0.9,0.9}{
            \makecell[{{p{\linewidth}}}]{
                \texttt{\tiny{[GM$|$GM]}}
                \texttt{yes} \\
            }
        }
    }
    & & \\ \\

    \theutterance \stepcounter{utterance}  
    & & & \multicolumn{2}{p{0.3\linewidth}}{
        \cellcolor[rgb]{0.9,0.9,0.9}{
            \makecell[{{p{\linewidth}}}]{
                \texttt{\tiny{[GM$|$GM]}}
                \texttt{game\_result = LOSE} \\
            }
        }
    }
    & & \\ \\

    \theutterance \stepcounter{utterance}  
    & & \multicolumn{4}{p{0.6\linewidth}}{
        \cellcolor[rgb]{0.95,0.95,0.95}{
            \makecell[{{p{\linewidth}}}]{
                \texttt{\tiny{[P1$\leftrightarrow$P]}}
                \texttt{Post{-}episode Reflexion (not sent to GM):} \\
\texttt{New plan: I failed because I initially answered "No" based on a misinterpretation of the counterfactual direction, then spent the rest of the response trying to self{-}correct and justify a "Yes" answer, which likely caused the system to reject the inconsistent output or fail to recognize the final correct state. The core logic was actually correct (\$69\textbackslash{}\% > 28\textbackslash{}\%\$), but the delivery was flawed due to the initial error. Next time, I will immediately start with the single word "Yes" as the very first token of the response, followed directly by the concise reasoning without any preliminary "No" or self{-}correction blocks. I will strictly adhere to the format: "Yes" followed by the explanation, ensuring no contradictory statements appear before the final conclusion.} \\
            }
        }
    }
    & \\ \\

\end{supertabular}
}

\end{document}
